% v11 Design: relocation, labels, guards and complexity. Replaces v10 "Relocation" and
% "State Invariant and Matching Limits". Numbers enter only through generated macros.
\subsection{Relocation}
Given $V_{t+1}$ and, when available, the previous text of $\pi$, relocation returns a new position and a label (Algorithm~\ref{alg:relocate}, Fig.~\ref{fig:cascade}). Tier~1 accepts an untouched file. Tier~2 searches $\qbar$ as a contiguous run of stripped non-blank lines in the same file and takes the occurrence nearest the old start. Tier~3 handles edits in place. It maps the old start line through a line diff \cite{pythonsequence} of the previous and new file to an expected start $p$, scores every window of height $k$ starting within $p \pm r$, $r = \max(10, 2k)$, by the similarity ratio of the joined stripped lines, and accepts the best window if the ratio reaches $\theta$ (deployed default $0.85$). Tier~4 searches $\qbar$ in the other files, in sorted path order. Tier~3 runs before tier~4 so that an edited block is recognised in place rather than matched to an identical copy elsewhere, and its window stays near the diff-predicted line to avoid distant look-alike code; Section~\ref{sec:design-results} measures both choices. A report anchor is matched in the same way against the sentences of the new report, with ratio threshold 0.8.

For normalized lines joined by a single space, let $u$ and $v$ be the resulting quotation and candidate strings. The fuzzy score is
\begin{equation}
\operatorname{ratio}(u,v)=\frac{2\sum_i m_i}{|u|+|v|},
\end{equation}
where $m_i$ are the matching-block lengths returned by Python's sequence matcher with its automatic junk heuristic disabled. Only full-height windows are considered. Starts are scanned in ascending order and the first maximum is kept; exact-run distance ties likewise select the first run. If the previous file text is unavailable, the predicted start defaults to the old start. These deterministic details decide ambiguous matches and are part of the evaluated implementation.

\begin{table}[t]
\caption{Relocation Labels: Definition, Limits and Effect on the Issue}\label{tab:labels}
\centering\footnotesize
\setlength{\tabcolsep}{3pt}
\begin{tabular}{@{}p{.17\columnwidth}p{.37\columnwidth}p{.24\columnwidth}p{.12\columnwidth}@{}}\toprule
Label & Holds when & Does \emph{not} establish & Issue \\ \midrule
\lbl{identical} & same path and same file digest (tier 1) & semantic review of the file & keeps state \\
\lbl{content-match} & normalized quotation $\qbar$ found at the old start (tier 2) & byte identity, equal indentation, equal behaviour & keeps state \\
\lbl{moved} & $\qbar$ found at another start or in another file (tiers 2, 4) & that the match is the original copy & keeps state \\
\lbl{changed} & best near-predicted window has ratio $\ge\theta$ (tier 3) & that the concern is affected & re-check \\
\lbl{removed} & no tier succeeds & that the code was deleted & re-check \\ \bottomrule
\end{tabular}
\tablenote{}{The deployed code emits a single identifier, \texttt{exact}, for both \lbl{identical} and \lbl{content-match}; this paper separates them because only the first is a byte-level property. Normalization strips each line's leading and trailing whitespace and skips blank lines.}
\end{table}

\subsection{Labels and Their Limits}
Table~\ref{tab:labels} defines the five labels. The word \emph{exact} in earlier versions of this work suggested more than the matcher establishes, so we use \lbl{content-match} for equality after normalization. Normalization is what lets an anchor survive blank-line and formatting edits, but in Python it also hides indentation, which carries meaning: a statement dedented out of a conditional keeps its stripped text and can keep its line number while its behaviour changes. A \lbl{content-match} or \lbl{moved} label is therefore textual evidence, not semantic evidence. A second limit concerns copies. When the original is edited or deleted and an identical copy exists elsewhere, tiers 2 and 4 can return the copy as \lbl{moved}; because \lbl{moved} keeps the issue's state, the change is not returned to the supervisor. Sections~\ref{sec:safety} and~\ref{sec:stress} measure both limits.

\subsection{Proposed Safety Guards}\label{sec:guards}
The failure analysis of Section~\ref{sec:safety} motivated three switches on the deployed cascade, and the compound-edit stress test a correction to one of them. They are evaluated alongside it under protocols whose text and code hashes were recorded before any held-out run (Section~\ref{sec:method}); the deployed cascade remains the system of record.
\begin{itemize}
\item \textbf{Indentation guard (IG).} When a normalized run is accepted at tier 2 or 4, the leading-whitespace widths of its non-blank lines (tabs expanded to eight columns) are compared with the quotation's, both relative to the first line and in absolute base indentation. For Python files any difference turns the label into \lbl{changed} with reason \emph{layout}; for other languages only a relative difference does. Placement is unchanged.
\item \textbf{Pre-existing-copy guard (AG).} At tier 2 a run that starts exactly at the diff-predicted line is preferred over the run nearest the old start. When a tier-2 or tier-4 run is accepted anywhere other than the predicted line, the guard counts occurrences of $\qbar$ in the old and new versions of the candidate file. A genuine move introduces the text, so only an increased count yields a clean \lbl{moved}; otherwise the placement is kept but re-check is requested with reason \emph{pre-existing copy}. The previous version's files are already loaded by carry-forward, so the guard needs no new input.
\item \textbf{Insertion fix (AGi).} AG as first specified flagged every clean insertion directly above an anchor: the pinned diff mapping sends the old start of such an anchor to the first inserted line, so the unmoved anchor never sits at its predicted start (Section~\ref{sec:compound}). AGi is AG with one change: a run at the line to which the diff maps the anchor's own first line also counts as the predicted start. It needs only the returned window and the two versions of the file, so it applies to the output of any locator (Section~\ref{sec:locator}).
\item \textbf{Variable-height window (VW).} Tier~3 also scores windows of heights $h\in[\max(1,k-d),\,k+d]$, $d=\max(1,\lceil k/2\rceil)$, so that statements that grow or shrink can be recovered with their new extent. Ties prefer the predicted start, then the original height.
\end{itemize}
None of the guards proves identity. Two copies introduced in the same revision defeat AG, a copy inserted just below the original can make the diff align the original's lines with the copy and so defeat AGi, and IG is a layout signal, not an equivalence test. We also tested a provenance rule (AG2) that accepts a window by where its lines came from in the diff instead of by copy counts. It removed more unnecessary requests but could not tell a move from an edited original with an identical copy elsewhere, and on the held-out duplicate suite it missed nearly as many re-checks as no guard at all (Section~\ref{sec:locator}). Content alone cannot decide which of two identical copies is the reviewed one; it can only decide when to ask.

\subsection{Keeping the Reviewed Text}\label{sec:identity}
When tier 3 accepts a \lbl{changed} window, carry-forward replaces the anchor's quotation with the matched text, and the next version is matched against that text. A wrong fuzzy match therefore becomes the reference for every later version, and a later \lbl{content-match} label then certifies only that the wrong code did not change again. We evaluate an identity rule (ID) that keeps the quotation the supervisor reviewed, $q_0$, unchanged beside the current one, $q_t$. Relocation still tracks $q_t$, but two checks use $q_0$. Every \lbl{changed} hop must also reach similarity 0.5 to $q_0$ (the calibrated threshold, Section~\ref{sec:design-results}); otherwise the transition is uncertain, the anchor is reported removed and its issue returns, so the anchor cannot drift step by step into unrelated code. And the request is computed against $q_0$: an issue returns for re-check whenever the located code differs from what the supervisor reviewed, however many versions ago that was. The interface would show $q_0$ beside $q_t$. A supervisor's confirmation sets a new $q_0$; the replay in Section~\ref{sec:prefixes} assumes no supervisor action. ID needs no particular locator and is evaluated over both the cascade and CodeMapper.

\subsection{Complexity}\label{sec:complexity}
Let the anchored file have $L$ lines, the version $N$ lines in $F$ files, the anchor $k$ lines and its joined quotation $q$ characters. Tier~1 compares stored digests in $O(1)$. Tier~2 normalizes the file and scans for runs in $O(L + ck)$, where $c$ counts lines equal to the first quoted line, so $O(Lk)$ in the worst case. Tier~3 first maps the old start through a line diff, quadratic in $L$ in the worst case for Python's sequence matcher \cite{pythonsequence}, and then scores $2r+1$ windows at $O(q^2)$ each, giving $O(L^2 + rq^2)$ with $r=\max(10,2k)$; VW multiplies the window term by $2d+1$. Tier~4 normalizes and scans every other file, $O(Nk)$. One anchor therefore costs $O(L^2 + rq^2 + Nk)$ in the worst case, and a review with $A$ anchors $O\!\left(A(L^2 + rq^2 + Nk)\right)$, because the implementation neither caches normalized files nor indexes lines across anchors. The $ANk$ term dominates for anchors whose code has disappeared from a large repository: every other file is scanned. A line index keyed by the first normalized line would reduce tier~4 to the cost of verifying candidate runs. Section~\ref{sec:scale} measures these costs on large repositories.
